% Tailored from reusable Faculty resume for Ohio State University R160519.
% Target: Assistant Professor, Optimization for Machine Learning (ISE/ECE).
\documentclass[9pt,letterpaper]{extarticle}
\usepackage[left=0.68in,right=0.68in,top=0.38in,bottom=0.38in]{geometry}
\usepackage[T1]{fontenc}
\input{glyphtounicode}
\pdfgentounicode=1
\usepackage[utf8]{inputenc}
\usepackage{lmodern}
\usepackage{microtype}
\usepackage{pifont}
\usepackage{xcolor}
\usepackage{enumitem}
\usepackage{titlesec}
\usepackage[hidelinks]{hyperref}
\usepackage{parskip,ragged2e,needspace}
\pagestyle{empty}
\setlength{\parindent}{0pt}
\setlength{\parskip}{1pt}
\hyphenpenalty=10000
\exhyphenpenalty=10000
\pretolerance=10000
\tolerance=2000
\emergencystretch=2em
\definecolor{accent}{HTML}{155B8A}
\titleformat{\section}{\large\bfseries\color{accent}}{}{0pt}{}[\vspace{0.5pt}\color{accent}\titlerule]
\titlespacing*{\section}{0pt}{2pt}{1pt}
\setlist[itemize]{leftmargin=1.15em,itemsep=0pt,topsep=0pt,parsep=0pt,partopsep=0pt}
\newcommand{\role}[4]{\Needspace{4\baselineskip}\textbf{#1}\hfill\textbf{#2}\\[-1pt]\textit{#3}\hfill\textit{#4}\par}
\newcommand{\project}[3]{\Needspace{5\baselineskip}\textbf{#1}\hfill\textit{#2 \textbar{} #3}\par}
\newcommand{\pub}[3]{\Needspace{3\baselineskip}\textbf{#1}\hfill\textit{#2}\par{\small #3\par}}
\hypersetup{
  pdftitle={Waseem Raza, PhD - Ohio State R160519 Optimization for Machine Learning},
  pdfauthor={Waseem Raza}
}

\begin{document}
\begin{center}
{\LARGE\bfseries Waseem Raza, PhD}\\[2pt]
{\large\bfseries Faculty Candidate | Optimization for Machine Learning \& AI-Enabled Networked Systems}\\[1pt]
{\small Robust Learning | Reinforcement Learning | Network Optimization | Scientific / Domain-Informed ML}\\[3pt]
{\footnotesize New Jersey, United States \quad
\ding{37}\enspace \href{tel:+17328039703}{+1 732-803-9703} \quad
\ding{41}\enspace \href{mailto:waseem.placements@gmail.com}{waseem.placements@gmail.com} \textbar{}
\href{mailto:waseemraza844@gmail.com}{waseemraza844@gmail.com}\\[1pt]
\href{https://www.linkedin.com/in/waseem-raza-phd-9a737858/}{LinkedIn} \textbar{}
\href{https://scholar.google.com/citations?user=d58Drr0AAAAJ&hl=en}{Google Scholar} \textbar{}
\href{https://orcid.org/0000-0002-5544-2637}{ORCID} \textbar{}
\href{https://github.com/WaseemRaza844}{GitHub} \textbar{}
\href{https://waseemraza844.github.io/portfolio/}{Portfolio}}
\end{center}
\vspace{-5pt}

\section{Academic Profile}
AI/ML researcher and engineering educator with a PhD in Electrical and Computer Engineering and research spanning optimization, reinforcement learning, robust machine learning, networked intelligent systems, 5G/6G, digital twins, and wearable intelligence. My work develops and rigorously evaluates data-driven methods for decision-making under limited data, distribution shift, heterogeneous environments, and operational uncertainty. Research experience combines analytical modeling, Python-based simulation, experimental testbeds, transfer learning, domain adaptation, generative modeling, statistical validation, and industry R\&D. Published work includes multi-parameter optimization and proactive reliability, robust propagation models for digital twins, intelligent outage management, user-centric RAN optimization, and AI-enabled sensing.

\section{Research Interests \& Methodological Fit}
\textbf{Optimization for ML \& AI-Enabled Decision Making:} multi-parameter optimization, sequential decision-making, reinforcement learning, network optimization, reliability-aware control, and data-driven decision support.\\
\textbf{Robust \& Trustworthy Learning:} domain adaptation, transfer learning, distribution-shift evaluation, sparse-data learning, robustness analysis, model validation, failure-mode assessment, and reliable AI for high-impact systems.\\
\textbf{Distributed / Communication-Constrained AI:} learning and optimization for wireless and edge systems, network-aware AI, mobility, radio access networks, measurement-driven modeling, and communication-constrained intelligent systems.\\
\textbf{Scientific \& Domain-Informed ML:} 3GPP-informed modeling, digital twins, propagation modeling, generative data augmentation, simulation-driven learning, statistical validation, and physics/domain-informed priors.\\
\textbf{Machine Learning Methods:} deep learning, reinforcement learning, GANs, autoencoders, classification, regression, anomaly detection, time-series analysis, cross-validation, and experimental evaluation.\\
\textbf{Research Tools:} Python, PyTorch, TensorFlow, scikit-learn, XGBoost, Pandas, NumPy, SQL, statistical analysis, reproducible experiments, and technical visualization.

\section{Research \& Professional Experience}
\role{Graduate Research Assistant}{09/2019 -- 07/2025}{AI4Networks, University of Oklahoma}{Tulsa, OK}
\begin{itemize}
\item Led wireless-AI research spanning radio propagation, digital twins, RAN optimization, and resilient networks.
\item Developed transfer learning, generative models, and reinforcement learning for sparse data.
\item Evaluated robustness under distribution shift and limited labels across 100+ experiments.
\item Built reusable Python simulations and an AI-enabled 5G testbed for repeatable studies.
\item Published and presented research internationally; mentored junior researchers in ML experiments and writing.
\end{itemize}

\role{Machine Learning Engineer}{02/2025 -- Present}{AT\&T, contractor via Innovatix Inc.}{Bedminster, NJ}
\begin{itemize}
\item Apply domain-adaptive ML to telecom operations, model validation, and network assurance.
\item Build domain-adaptive neural models to improve generalization across operational environments.
\item Develop cross-domain validation workflows with telecom, software, and public-safety teams.
\end{itemize}

\role{Advanced-Degrees Intern}{06/2024 -- 08/2024}{AT\&T Labs}{Austin, TX}
\begin{itemize}
\item Analyzed 500K+ LTE/5G-NR logs, using ML to reduce diagnostic time by 20\%.
\item Improved 5G NR-NTN cell-selection algorithms through experimental analysis and validation.
\end{itemize}

\role{Advanced Wireless Technology Development Intern}{09/2022 -- 01/2023}{Futurewei Technologies}{Schaumburg, IL}
\begin{itemize}
\item Improved wireless CSI compression by 25\% through transfer learning and experiments.
\item Completed 50+ experiments assessing robustness and generalization in wireless ML.
\end{itemize}

\section{Education}
\textbf{PhD, Electrical and Computer Engineering}\hfill\textit{09/2019 -- 07/2025}\\
University of Oklahoma, Norman, Oklahoma, USA\\
\textbf{MSc, Telecommunication Engineering}\hfill\textit{08/2014 -- 08/2016}\\
University of Engineering and Technology, Taxila, Pakistan\\
\textbf{BSc, Telecommunication Engineering}\hfill\textit{11/2010 -- 06/2014}\\
University of Engineering and Technology, Taxila, Pakistan

\section{Selected Publications}
\pub{\href{https://doi.org/10.3390/s25247651}{On Multi-Parameter Optimization and Proactive Reliability in 5G and Beyond Cellular Networks}}{2025}{Sensors --- data-driven multi-parameter optimization, proactive reliability, and decision support for future cellular systems.}
\pub{\href{https://doi.org/10.1109/PIMRC59610.2024.10817403}{An AI-Driven Framework for Enhancing Resilience in Propagation Models to Enable Digital Twin}}{2024}{IEEE PIMRC --- robust data-driven propagation modeling, generative learning, and resilience under scarce measured data and distribution shift.}
\pub{Data-Driven Intelligent Outage Management for High Shadowing Environments in 5G\&B Networks}{2024}{IEEE SmartNets --- machine-learning-based outage detection, intelligent management, and reliability in challenging radio environments.}
\pub{Towards Positioning Error Impact Characterization and Minimization in User-Centric RAN}{2022}{IEEE WCNC --- data-driven characterization and optimization for user-centric radio access networks.}
\pub{\href{https://arxiv.org/abs/2203.15899}{Toward a Hybrid RF/Optical Lunar Communication System (LunarComm)}}{2022}{IEEE Network --- architecture and design tradeoffs for resilient, high-data-rate hybrid networked systems.}
\pub{\href{https://doi.org/10.1016/j.aeue.2020.153244}{Unification of RF Energy Harvesting Schemes under Mixed Rayleigh--Rician Fading Channels}}{2020}{AEÜ -- International Journal of Electronics and Communications --- analytical modeling of energy-harvesting cooperative-relaying systems under mixed fading conditions.}
\pub{Real-World Medication Adherence Monitoring Using Hybrid Neural Networks and Smartwatches}{2026}{HealthINF (Accepted) --- deep learning, multimodal wearable sensing, cross-user evaluation, and reliable human-centered AI.}

\section{Selected Research Projects}
\project{Multi-Parameter Optimization, Reliability \& Self-Healing for 5G/6G}{Research Program}{2020--2025}
Developed data-driven methods for outage management, proactive reliability, network optimization, and autonomous decision support in cellular systems. The work connects learning-based models with sequential network decisions and reliability objectives, including robustness across changing radio environments.\\[-1pt]
{\small\textit{Methods \& tools:} reinforcement learning, network optimization, reliability analysis, outage management, deep learning, statistical modeling, Python, model validation, 5G/6G systems.}

\project{Propagation Models \& Wireless Digital Twins}{Research Project}{2024}
Developed robust propagation models for wireless digital twins using domain-informed generative learning and scarce measured data. Combined 3GPP-informed pretraining, GAN-based augmentation, domain adaptation, and statistical validation to improve resilience under distribution shift.\\[-1pt]
{\small\textit{Methods \& tools:} propagation modeling, digital twins, GANs, domain adaptation, 3GPP models, synthetic data, statistical validation, robustness evaluation, Python.}

\project{SyntheticNET: 3GPP-Compliant Network Simulation}{Research Project}{2020--2024}
Developed a Python-based, 3GPP-aligned simulation environment for realistic mobility, propagation, and cellular-network experimentation when measured data is sparse or costly. Enabled controlled data generation and repeatable studies in optimization, digital twins, and resilient wireless AI.\\[-1pt]
{\small\textit{Methods \& tools:} Python, 3GPP, wireless simulation, mobility modeling, propagation modeling, synthetic data, experiment design, network analytics.}

\project{TurboRAN Experimental Platform}{Research Infrastructure}{2020--2024}
Built and used a multi-band, multi-tier experimental platform for evaluating AI-enabled radio-access-network methods under controlled conditions. Supported repeatable network measurements, testbed experimentation, RAN optimization studies, and validation of emerging algorithms.\\[-1pt]
{\small\textit{Methods \& tools:} 5G testbeds, RAN, network measurement, experiment design, wireless prototyping, performance evaluation, optimization.}

\Needspace{3\baselineskip}
\section{Academic Appointments, Teaching \& Mentoring}
\role{Lecturer}{Prior Academic Appointment}{Engineering and Technology Programs}{}
\begin{itemize}
\item Taught undergraduate communications and programming through applied engineering problems.
\item Mentored student researchers in ML experiments, reproducibility, and research writing.
\end{itemize}

\role{Research Associate, ComSens Lab}{06/2016 -- 12/2017}{COMSATS University Islamabad}{Islamabad, Pakistan}
\begin{itemize}
\item Conducted wireless-communications research using simulation and quantitative modeling.
\item Mentored BSc and MSc thesis projects in experiment design and research presentation.
\item Reviewed technical reports and supported proposals, seminars, and research coordination.
\end{itemize}

\role{Full-Time Research Scholar}{08/2014 -- 02/2016}{UET Taxila}{Taxila, Pakistan}
\begin{itemize}
\item Researched communication systems while completing an MSc in Telecommunication Engineering.
\item Assisted communications laboratories, assignments, grading, and examination preparation.
\end{itemize}

\Needspace{3\baselineskip}
\section{Teaching Areas \& ISE/ECE Contribution}
\begin{itemize}
\item Prepared to teach programming, machine learning, signals, communications, and labs.
\item Graduate interests: optimization, robust ML, reinforcement learning, and networked systems.
\item Use project-based learning for reproducible implementation and quantitative evaluation.
\item Mentor ISE/ECE projects in optimization, wireless AI, and networked intelligent systems.
\end{itemize}

\Needspace{3\baselineskip}
\section{Research Trajectory for Optimization for Machine Learning}
My research will advance optimization and learning for reliable networked intelligent systems through: (1) robust ML under distribution shift, uncertainty, and limited labels; (2) sequential reinforcement-learning methods for autonomous network control; (3) distributed, communication-efficient learning for wireless and edge systems; and (4) trustworthy AI with rigorous reliability, robustness, and interpretability evaluation. This agenda extends my work in network optimization, reinforcement learning, domain adaptation, digital twins, outage management, and experimental wireless AI.

\Needspace{3\baselineskip}
\section{Collaboration, Research Translation \& Academic Service}
\begin{itemize}
\item Collaborate with academic and industry teams on measurable AI research and validation.
\item Bridge academic and industrial R\&D through reproducible experiments, peer review, and mentoring.
\end{itemize}

\end{document}
